% Source: https://artofproblemsolving.com/community/c3h3579385_1434th_post
% Stable identity: aops-topic-3579385-q16
For a prime $p$, write $v_p(i)$ for its exponent in $i$. Since $v_p(f(i))=2\lfloor\dfrac{v_p(i)}{2}\rfloor$, the exponent of $p$ in $Q_n$ is\[E_p=\sum_{j\ge1}\left(\left\lfloor\dfrac{n}{p^{2j-1}}\right\rfloor-2\left\lfloor\dfrac{n}{p^{2j}}\right\rfloor\right).\]For $p\ge3$, each parenthesis is nonnegative because its first floor is at least $p$ times the second floor. For $p=2$, each parenthesis is exactly the binary digit of $n$ in position $2j-1$, numbering the units position zero. Thus all prime exponents are nonnegative, so $Q_n$ is always an integer. Since $n<2^{10}$, its exponent of $2$ is the number of ones among positions $1,3,5,7,9$. It equals two precisely when two of these five positions contain ones. The other five binary digits are arbitrary. Therefore there are $\binom52\cdot2^5=\boxed{320}$ choices. The excluded value $n=0$ has no selected ones, so is not counted.
